\section{Introduction}
\label{sec:intro}

A market maker earns the spread by resting limit orders on both sides of the
book. Whether a given order earns anything depends on a single event: whether an
aggressor reaches it before its owner withdraws it. In a backtest that event can
never be observed, because the order being tested was never in the market. The
simulator must therefore decide, by some rule, whether the order \emph{would}
have been filled. Candle-based platforms fill a limit order once the bar's
extreme touches its price; event-driven simulators wait until enough volume has
traded at the price; stochastic-control models of market making replace the
question with a fill intensity that decays with the distance from the mid-price.

These rules are not innocuous approximations. They determine the fill rate,
hence the inventory path, the realised spread, the adverse selection a strategy
suffers and, ultimately, which of two strategies looks better. Yet they are
almost never validated, for the simple reason that the data needed to validate
them---the outcome of each individual resting order together with the exact
queue in front of it---has not been publicly available.

The order-level dataset of \citet{albers2026openbook} changes this. Hyperliquid,
a perpetual-futures exchange whose central limit order book runs on its own
blockchain, publishes every order of every participant together with its fate.
The dataset therefore contains, for millions of real passive orders, exactly the
quantity a backtest has to guess: whether the order was filled, and when. This
lets us ask, and answer, a question that has so far been inaccessible:
\emph{how much does the unobservable fill assumption change the conclusions of a
market-making backtest?}

We answer in three steps, each resting on the one before.

\paragraph{Step one: how wrong the rules are.}
We rebuild the queue in front of every real passive order and compare each fill
rule with what actually happened (Section~\ref{sec:fillmodels}). The touch rule
overstates fills five- to seven-fold and the candle rule of bar-based backtests
by far more; the volume-through rule understates them two- to three-fold. Both
errors have one cause: more than 94\% of the volume queued ahead of an order
leaves by cancellation, not by trading. A rule can also rank orders almost
perfectly and still be worse, as a yes/no answer, than assuming that nothing is
ever filled.

\paragraph{Step two: what it takes to be right.}
A rule that follows the queue comes within a few percent of the truth, and its
whole remaining error is the exchange's processing order within a block:
cancellations are taken before aggressors, which rescues 3--4\% of orders that a
trade would otherwise have filled. With that rule the model reproduces the
observed fills to within 0.3\% (Section~\ref{sec:block}). We then ask how far one
can get without order-level data (Section~\ref{sec:learned}), and how much of an
order's fate can be known when it is placed at all: fill times can be ranked
with a concordance of at most 0.92--0.93, yet placement-time information explains
little of which orders are actually filled (Section~\ref{sec:knowable}).

\paragraph{Step three: what the error does to a backtest.}
We backtest 48 variants of twelve families of quoting strategies under each fill
rule, changing nothing else (Section~\ref{sec:backtest}). Under the validated
rule 47 of them lose money on every instrument. The touch rule flatters 75--92\%
of them; the exponential clock turns 83--90\% of the losses into profits and
ranks the strategies in an order unrelated to the true one. Neither error is a
uniform shift: each rewards a particular style of quoting.

\paragraph{Reproducibility.}
Every number in this paper is printed by a script in the accompanying
repository and its reference output is stored alongside the code
(Appendix~\ref{app:repro}).
